%%% Hlavní soubor. Zde se definují základní parametry a odkazuje se na ostatní části. %%%

%% Verze pro jednostranný tisk:
% Okraje: levý 40mm, pravý 25mm, horní a dolní 25mm
% (ale pozor, LaTeX si sám přidává 1in)
\documentclass[12pt,a4paper]{report}
\setlength\textwidth{145mm}
\setlength\textheight{247mm}
\setlength\oddsidemargin{15mm}
\setlength\evensidemargin{15mm}
\setlength\topmargin{0mm}
\setlength\headsep{0mm}
\setlength\headheight{0mm}
% \openright zařídí, aby následující text začínal na pravé straně knihy
\let\openright=\clearpage

%% Pokud tiskneme oboustranně:
% \documentclass[12pt,a4paper,twoside,openright]{report}
% \setlength\textwidth{145mm}
% \setlength\textheight{247mm}
% \setlength\oddsidemargin{14.2mm}
% \setlength\evensidemargin{0mm}
% \setlength\topmargin{0mm}
% \setlength\headsep{0mm}
% \setlength\headheight{0mm}
% \let\openright=\cleardoublepage

% Přepneme na českou sazbu
\usepackage[slovak]{babel}
\usepackage[IL2]{fontenc}

%% Použité kódování znaků: obvykle latin2, cp1250 nebo utf8:
\usepackage[utf8]{inputenc}

%%% Další užitečné balíčky (jsou součástí běžných distribucí LaTeXu)
\usepackage{amsmath}        % rozšíření pro sazbu matematiky
\usepackage{amsfonts}       % matematické fonty
\usepackage{amsthm}         % sazba vět, definic apod.
\usepackage{bbding}         % balíček s nejrůznějšími symboly
			    % (čtverečky, hvězdičky, tužtičky, nůžtičky, ...)
\usepackage{bm}             % tučné symboly (příkaz \bm)
\usepackage{graphicx}       % vkládání obrázků
\usepackage{fancyvrb}       % vylepšené prostředí pro strojové písmo
\usepackage{indentfirst}    % zavede odsazení 1. odstavce kapitoly
\usepackage{natbib}         % zajištuje možnost odkazovat na literaturu
			    % stylem AUTOR (ROK), resp. AUTOR [ČÍSLO]
\usepackage[nottoc]{tocbibind} % zajistí přidání seznamu literatury,
                            % obrázků a tabulek do obsahu
\usepackage{icomma}         % inteligetní čárka v matematickém módu
\usepackage{dcolumn}        % lepší zarovnání sloupců v tabulkách
\usepackage{booktabs}       % lepší vodorovné linky v tabulkách
\usepackage{paralist}       % lepší enumerate a itemize
\usepackage[usenames]{xcolor}  % barevná sazba
\usepackage{float}

\usepackage{subcaption}
\usepackage[noend]{algpseudocode}
\usepackage{algorithm}
\usepackage{cases}
\usepackage{tikz}
\usepackage{amssymb}
\usepackage{xspace}
\usetikzlibrary{positioning, shapes, trees, graphs} % RNA trees
\usetikzlibrary{decorations.pathmorphing} % noisy shapes
\usetikzlibrary{fit}					% fitting shapes to coordinates
\usetikzlibrary{backgrounds}	% drawing the background after the foreground

\usepackage[percent]{overpic}


%%% Údaje o práci

% Název práce v jazyce práce (přesně podle zadání)
\def\NazevPrace{Analýza datových toků ve databázových systémech}

% Název práce v angličtině
\def\NazevPraceEN{Analyzing Data Lineage in Database Frameworks}

% Jméno autora
\def\AutorPrace{Richard Eliáš}

% Rok odevzdání
\def\RokOdevzdani{2019}

% Název katedry nebo ústavu, kde byla práce oficiálně zadána
% (dle Organizační struktury MFF UK, případně plný název pracoviště mimo MFF)
\def\Katedra{Katedra distribuovaných a spolehlivých systémů}
\def\KatedraEN{Department of Distributed and Dependable Systems}

% Jedná se o katedru (department) nebo o ústav (institute)?
\def\TypPracoviste{Katedra}
\def\TypPracovisteEN{Department}

% Vedoucí práce: Jméno a příjmení s~tituly
\def\Vedouci{RNDr. Pavel Parízek, Ph.D.}

% Pracoviště vedoucího (opět dle Organizační struktury MFF)
\def\KatedraVedouciho{Katedra distribuovaných a spolehlivých systémů}
\def\KatedraVedoucihoEN{Department of Distributed and Dependable Systems}

% Studijní program a obor
\def\StudijniProgram{Informatika}
\def\StudijniObor{Umělá inteligence}

% Nepovinné poděkování (vedoucímu práce, konzultantovi, tomu, kdo
% zapůjčil software, literaturu apod.)
\def\Podekovani{%
TODO Podakovanie
}

% Abstrakt (doporučený rozsah cca 80-200 slov; nejedná se o zadání práce)
\def\Abstrakt{%
TODO Abstrakt
}
\def\AbstraktEN{%
TODO Abstrakt
}

% 3 až 5 klíčových slov (doporučeno), každé uzavřeno ve složených závorkách
\def\KlicovaSlova{%
  TODO Klucove slova
}
\def\KlicovaSlovaEN{%
  TODO Klucove slova
}

%% Balíček hyperref, kterým jdou vyrábět klikací odkazy v PDF,
%% ale hlavně ho používáme k uložení metadat do PDF (včetně obsahu).
\usepackage[pdftex,unicode]{hyperref}   % Musí být za všemi ostatními balíčky
\hypersetup{breaklinks=true}
\hypersetup{pdftitle={\NazevPrace}}
\hypersetup{pdfauthor={\AutorPrace}}
\hypersetup{pdfkeywords=\KlicovaSlova}
\hypersetup{urlcolor=blue}

%% Definice různých užitečných maker (viz popis uvnitř souboru)
%%% Tento soubor obsahuje definice různých užitečných maker a prostředí %%%
%%% Další makra připisujte sem, ať nepřekáží v ostatních souborech.     %%%

%%% Drobné úpravy stylu

% Tato makra přesvědčují mírně ošklivým trikem LaTeX, aby hlavičky kapitol
% sázel příčetněji a nevynechával nad nimi spoustu místa. Směle ignorujte.
\makeatletter
\def\@makechapterhead#1{
  {\parindent \z@ \raggedright \normalfont
   \Huge\bfseries \thechapter. #1
   \par\nobreak
   \vskip 20\p@
}}
\def\@makeschapterhead#1{
  {\parindent \z@ \raggedright \normalfont
   \Huge\bfseries #1
   \par\nobreak
   \vskip 20\p@
}}
\makeatother

% Toto makro definuje kapitolu, která není očíslovaná, ale je uvedena v obsahu.
\def\chapwithtoc#1{
\chapter*{#1}
\addcontentsline{toc}{chapter}{#1}
}

% Trochu volnější nastavení dělení slov, než je default.
\lefthyphenmin=2
\righthyphenmin=2

% Zapne černé "slimáky" na koncích řádků, které přetekly, abychom si
% jich lépe všimli.
\overfullrule=1mm

%%% Makra pro definice, věty, tvrzení, příklady, ... (vyžaduje baliček amsthm)

\theoremstyle{plain}
\newtheorem{veta}{Veta}
\newtheorem{lemma}[veta]{Lemma}
\newtheorem{tvrz}[veta]{Tvdenie}

\theoremstyle{plain}
\newtheorem{definice}{Definícia}

\theoremstyle{remark}
\newtheorem*{dusl}{Dôsledok}
\newtheorem*{pozn}{Poznámka}
\newtheorem*{prikl}{Príklad}

%%% Prostředí pro důkazy

\newenvironment{dukaz}{
  \par\medskip\noindent
  \textit{Dôkaz}.
}{
\newline
\rightline{$\square$}  % nebo \SquareCastShadowBottomRight z balíčku bbding
}

%%% Prostředí pro sazbu kódu, případně vstupu/výstupu počítačových
%%% programů. (Vyžaduje balíček fancyvrb -- fancy verbatim.)

\DefineVerbatimEnvironment{code}{Verbatim}{fontsize=\small, frame=single}

%%% Prostor reálných, resp. přirozených čísel
\newcommand{\R}{\mathbb{R}}
\newcommand{\N}{\mathbb{N}}

%%% Užitečné operátory pro statistiku a pravděpodobnost
\DeclareMathOperator{\pr}{\textsf{P}}
\DeclareMathOperator{\E}{\textsf{E}\,}
\DeclareMathOperator{\var}{\textrm{var}}
\DeclareMathOperator{\sd}{\textrm{sd}}

%%% Příkaz pro transpozici vektoru/matice
\newcommand{\T}[1]{#1^\top}

%%% Vychytávky pro matematiku
\newcommand{\goto}{\rightarrow}
\newcommand{\gotop}{\stackrel{P}{\longrightarrow}}
\newcommand{\maon}[1]{o(n^{#1})}
\newcommand{\abs}[1]{\left|{#1}\right|}
\newcommand{\dint}{\int_0^\tau\!\!\int_0^\tau}
\newcommand{\isqr}[1]{\frac{1}{\sqrt{#1}}}

%%% Vychytávky pro tabulky
\newcommand{\pulrad}[1]{\raisebox{1.5ex}[0pt]{#1}}
\newcommand{\mc}[1]{\multicolumn{1}{c}{#1}}

\renewcommand{\O}[1]{\sloppy\mbox{\ensuremath{\mathcal{O}}(#1)}}

\newcommand{\scale}{0.6}
\newcommand{\subtree}[1]{\begin{tikzpicture}[
          on grid,
          font=\tiny,
          scale = \scale,
          node distance = 0.6 cm,
          sibling distance = 0.8 cm,
          text width = 0.2cm,
          level distance = 1.1 cm,
          colored/.style = {color = red},
          normal/.style = {black = red, draw, opacity = 100},
          invisible/.style={opacity = 0},
          cross/.style={path picture={ 
                \draw
                (path picture bounding box.south east) -- (path picture bounding box.north west) (path picture bounding box.south west) -- (path picture bounding box.north east);
            }},
          every node/.style = {scale = \scale, circle, draw, align = center},
          tree/.style = {draw = none, fill = none}]#1\end{tikzpicture}}

\definecolor{green}{rgb}{0,0.6,0}
\definecolor{gray}{rgb}{0.5,0.5,0.5}
\definecolor{yellow}{rgb}{0.73,0.71,0.16}
\definecolor{pink}{rgb}{0.58,0,0.82}
\definecolor{orange}{rgb}{0.8,0.47,0.19}

\lstdefinelanguage{Java}{
  morekeywords={
    return,try,catch,finally,
    static,default,final,volatile,new,
    null,void,int,long,boolean,short,double,float,
    interface,class,extends,implements,
    if,else,for,while,do,throw,throws,
    public,protected,private},
  numbers=left,
  numbersep=3mm,
  numberstyle=\tiny\color{gray},
  frame=tb,
  %aboveskip=3mm,
  %belowskip=3mm,
  showstringspaces=false,
  columns=flexible,
  basicstyle={\footnotesize\ttfamily},
  showlines=false,
  keywordstyle=\color{blue},
  commentstyle=\color{green},
  morecomment=[l]{//},
  morecomment=[s]{/*}{*/},
  stringstyle=\color{pink},
  morestring=[b]",
  morestring=[d]’,
  moredelim=[is][\textcolor{yellow}]{@@}{@@},
  breaklines=true,
  breakatwhitespace=true,
  escapeinside={/*}{*/},
  captionpos=b,
}

\lstdefinelanguage{XML}{
  %language=HTML,
  numbers=left,
  numbersep=3mm,
  numberstyle=\tiny\color{gray},
  frame=tb,
  %aboveskip=3mm,
  %belowskip=3mm,
  showstringspaces=false,
  columns=flexible,
  basicstyle={\footnotesize\ttfamily},
  showlines=false,
  keywordstyle=\color{blue},
  commentstyle=\color{green},
  morecomment=[s]{<!--}{-->},
  stringstyle=\color{pink},
  moredelim=[is][\textcolor{blue}]{@}{@},
  moredelim=[is][\textcolor{orange}]{@@}{@@},
  morestring=[b]",
  morestring=[d]’,
  breaklines=true,
  breakatwhitespace=true,
  escapeinside={/*}{*/},
  captionpos=b,
}

\makeatletter
\AtBeginDocument{%
  \let\c@table\c@lstlisting
  \let\thetable\thelstlisting

  \let\c@figure\c@lstlisting
  \let\thefigure\thelstlisting
  \let\ftype@lstlisting\ftype@figure
}
\makeatother

\newcommand\TODO[1]{{\color{red}TODO #1}}
\newcommand{\uvodzovky}[1]{``#1''}



%% Titulní strana a různé povinné informační strany
\begin{document}
%%% Title page of the thesis and other mandatory pages

%%% Title page of the thesis

\pagestyle{empty}
\hypersetup{pageanchor=false}
\begin{center}

\centerline{\mbox{\includegraphics[width=166mm]{img/logo-en.pdf}}}

\vspace{-8mm}
\vfill

{\bf\Large MASTER THESIS}

\vfill

{\LARGE\AutorPrace}

\vspace{15mm}

{\LARGE\bfseries\NazevPrace}

\vfill

\Katedra

\vfill

\begin{tabular}{rl}
Supervisor of the master thesis: & \Vedouci \\
\noalign{\vspace{2mm}}
Study programme: & \StudijniProgram \\
\noalign{\vspace{2mm}}
Study branch: & \StudijniObor \\
\end{tabular}

\vfill

% Zde doplňte rok
Prague \RokOdevzdani

\end{center}

\newpage

%%% Here should be a bound sheet included -- a signed copy of the "master
%%% thesis assignment". This assignment is NOT a part of the electronic
%%% version of the thesis. DO NOT SCAN.

%%% A page with a solemn declaration to the master thesis

\openright
\hypersetup{pageanchor=true}
\pagestyle{plain}
\pagenumbering{roman}
\vglue 0pt plus 1fill

I~hereby declare that I~have authored this thesis independently,
and that all sources used are declared in accordance with the
\uvodzovky{Metodický pokyn o~etické přípravě vysokoškolských závěrečných prací}.

\vspace{2mm}
I~acknowledge that my thesis (work) is subject to the rights and obligations
arising from Act No. 121/2000 Coll., on Copyright and Rights Related to Copyright
and on Amendments to Certain Laws (the Copyright Act), as amended,
(hereinafter as the "Copyright Act"), in particular §~35, and §~60 of the Copyright Act
governing the school work.

\vspace{2mm}
With respect to the computer programs that are part of my thesis (work)
and with respect to all documentation related to the computer programs ("software"),
I~hereby grant the so-called MIT License.

\vspace{2mm}
The MIT License represents a~license to use the software free of charge.
I~grant this license to every person interested in using the software.
Each person is entitled to obtain a~copy of the software (including
the related documentation) without any limitation, and may, without limitation,
use, copy, modify, merge, publish, distribute, sublicense and~/~or sell
copies of the software, and allow any person to whom the software is further
provided to exercise the aforementioned rights. Ways of using the software or the extent
of this use are not limited in any way.

\vspace{2mm}
The person interested in using the software is obliged to attach the text of the license terms as follows:

\vspace{3mm}
\noindent
Copyright (c) \RokOdevzdani~\AutorPrace \\
Permission is hereby granted, free of charge, to any person \\
obtaining a copy of this software and associated documentation \\
files (the \uvodzovky{Software}), to deal in the Software without \\
restriction, including without limitation the rights to use, \\
copy, modify, merge, publish, distribute, sublicense, and/or sell \\
copies of the Software, and to permit persons to whom the \\
Software is furnished to do so, subject to the following conditions: \\
The above copyright notice and this permission notice shall be \\
included in all copies or substantial portions of the Software.

{
\footnotesize
\vspace{3mm}
\noindent
THE SOFTWARE IS PROVIDED \uvodzovky{AS IS}, WITHOUT WARRANTY OF ANY KIND, \\
EXPRESS OR IMPLIED, INCLUDING BUT NOT LIMITED TO THE WARRANTIES \\
OF MERCHANTABILITY, FITNESS FOR A PARTICULAR PURPOSE AND \\
NONINFRINGEMENT. IN NO EVENT SHALL THE AUTHORS OR COPYRIGHT \\
HOLDERS BE LIABLE FOR ANY CLAIM, DAMAGES OR OTHER LIABILITY, \\
WHETHER IN AN ACTION OF CONTRACT, TORT OR OTHERWISE, ARISING \\
FROM, OUT OF OR IN CONNECTION WITH THE SOFTWARE OR THE USE OR \\
OTHER DEALINGS IN THE SOFTWARE.
}

\vspace{15mm}

\hbox{\hbox to 0.5\hsize{%
In $\ldots\ldots\ldots$ date $\ldots\ldots\ldots$
\hss}}

\vspace{20mm}
\newpage

%%% Podekovani

\openright

\Podakovanie

\newpage

%%% Mandatory information page of the thesis

\openright

\vbox to 0.5\vsize{
\setlength\parindent{0mm}
\setlength\parskip{5mm}

Title:
\NazevPrace

Author:
\AutorPrace

\TypPracoviste:
\Katedra

Supervisor:
\Vedouci, \KatedraVedouciho

Abstract:
\Abstrakt

Keywords:
\KlicovaSlova

\vss}

\newpage

\openright
\pagestyle{plain}
\pagenumbering{arabic}
\setcounter{page}{1}



%%% Strana s automaticky generovaným obsahem bakalářské práce

\tableofcontents

%%% Jednotlivé kapitoly práce jsou pro přehlednost uloženy v samostatných souborech

\chapter{Introduction}

For modern data processing applications, an important aspect
is a data lineage.
Data lineage, or provenance, describes where data came from,
how it was derived and where the results are stored.

Lineage can be useful in many domains.
Molecular biology databases, which mostly store copied data,
can use lineage to verify the copied data by tracking the original
sources as was demonstrated by \citet{DataLineageInBioInformatics}.
Probabilistic databases can exploit lineage for confidence
computation as was stated by \citet{DataLineageInProbabilisticDatabases}.
The important usage of the data lineage is also in machine learning,
where for the purpose of a verification and validation can be important
to know data sources of used data sets.
Data lineage is also important in Business Intelligence (BI)
when achieving full regulatory compliance and improving data governance.

In all these areas, Java applications together with data manipulation frameworks
are often used to access data in databases, files or in network storages.
The tool that is able to identify the sources and sinks of a data
and create data lineage of an application can be very useful.

Up to now, we are aware of the only tool for creating data lineage of
Java applications - the Symbolic analysis library
presented in \citet{ParizekHybridAnalysis}.
The library can create data lineage of an analysed application when only
Java I/O and JDBC APIs are used in addition to the core
Java libraries (collections, etc.).



\section{MANTA Flow}

An example of a data lineage analysis tool is the \citet{MantaFlow}.
The MANTA Flow helps enterprises to get end-to-end
data lineage including custom SQL code. That allows customers to fulfill
compliance regulations or improve data governance.

It extracts and analyses metadata from report definitions,
custom SQL code, and extract-transform-load (ETL) workflows, to create data flow graphs
which span multiple systems and a range of technologies.
Lineages are computed based on analysis of actual code.
All entities detected by MANTA Flow, such as database tables, columns or procedures,
are visualized to help users utilize this information.

Systems, for which MANTA Flow is used to analyse their data lineage,
often use Java applications and data processing frameworks to work with data.
MANTA Flow integrates the Symbolic analysis library to create data lineage of
Java applications to cover whole system, not only its database parts.
However, the support for computing data lineage of complex data processing frameworks
is missing.




\section{Goals}

In our work, we \textit{propose architecture changes} to Symbolic analysis library
to be easily extensible by plugins that can add support for new data processing frameworks
to identify its data sources and sinks.

We also \textit{implement Symbolic analysis library plugins} for few selected frameworks
(MyBatis, Spring JDBC and Apache Kafka).
Each framework has different approach for accessing the data.
Thereby we want to demonstrate that such plugins can be used to
add support for data lineage analysis of other Java frameworks.
The feature to easily extend the Symbolic analysis library to add support for a new frameworks
is very important, as new frameworks are being developed today.




\section{Structure of the Work}

The rest of the thesis has the following structure.

In chapter \ref{chapter:frameworks} we give overview of popular Java frameworks
for data processing.

In chapter \ref{chapter:analysis} we introduce libraries that are used
for static analysis of Java programs to create data lineage.

In chapter \ref{chapter:implementation} we present design of the \ToolName tool for
data lineage of frameworks and selected implementation details.

Chapter \ref{chapter:program} contains user documentation for the \ToolName tool.

Chapter \ref{chapter:results} present our results for selected frameworks,
limitations of our solution and plans for the future work.


\chapter*{Záver}
\addcontentsline{toc}{chapter}{Záver}

TODO Nejaky zaver


%%% Seznam použité literatury
%%% Seznam použité literatury (bibliografie)
%%%
%%% Pro vytváření bibliografie používáme bibTeX. Ten zpracovává
%%% citace v textu (např. makro \cite{...}) a vyhledává k nim literaturu
%%% v souboru literatura.bib.
%%%
%%% Příkaz \bibliographystyle určuje, jakým stylem budou citovány odkazy
%%% v textu. V závorce je název zvoleného souboru .bst. Styly plainnat
%%% a unsrt jsou standardní součástí latexových distribucí. Styl czplainnat
%%% je dodáván s touto šablonou a bibTeX ho hledá v aktuálním adresáři.

\bibliographystyle{plainnat}    %% Autor (rok) s anglickými spojkami
% \bibliographystyle{unsrt}       %% [číslo]

\renewcommand{\bibname}{References}

%%% Vytvoření seznamu literatury. Pozor, pokud jste necitovali ani jednu
%%% položku, seznam se automaticky vynechá.

\bibliography{literatura}

%%% Kdybyste chtěli bibliografii vytvářet ručně (bez bibTeXu), lze to udělat
%%% následovně. V takovém případě se řiďte normou ISO 690 a zvyklostmi v oboru.

% \begin{thebibliography}{99}
%
% \bibitem{lamport94}
%   {\sc Lamport,} Leslie.
%   \emph{\LaTeX: A Document Preparation System}.
%   2. vydání.
%   Massachusetts: Addison Wesley, 1994.
%   ISBN 0-201-52983-1.
%
% \end{thebibliography}


%%% Obrázky v bakalářské práci
%%% (pokud jich je malé množství, obvykle není třeba seznam uvádět)
\listoffigures

%%% Tabulky v bakalářské práci (opět nemusí být nutné uvádět)
%%% U matematických prací může být lepší přemístit seznam tabulek na začátek práce.
%\listoftables

%%% Přílohy k bakalářské práci, existují-li. Každá příloha musí být alespoň jednou
%%% odkazována z vlastního textu práce. Přílohy se číslují.
%%%
%%% Do tištěné verze se spíše hodí přílohy, které lze číst a prohlížet (dodatečné
%%% tabulky a grafy, různé textové doplňky, ukázky výstupů z počítačových programů,
%%% apod.). Do elektronické verze se hodí přílohy, které budou spíše používány
%%% v elektronické podobě než čteny (zdrojové kódy programů, datové soubory,
%%% interaktivní grafy apod.). Elektronické přílohy se nahrávají do SISu a lze
%%% je také do práce vložit na CD/DVD.
%\chapwithtoc{Prílohy}

\openright
\end{document}
